\documentclass[11pt]{article}
\usepackage[utf8]{inputenc}
\usepackage[T1]{fontenc}
\usepackage[margin=1in]{geometry}
\usepackage[hidelinks]{hyperref}
\usepackage{parskip}
\setlength{\parskip}{0.5em}
\pagestyle{empty}

\begin{document}

\noindent\textbf{\large Cover Letter --- \textit{Fractal and Fractional}}

\vspace{0.1em}

\noindent Dear Editors,

\vspace{0.15em}

\noindent I am pleased to submit our manuscript, \textbf{``Spectral Isomorphism between Renormalization Flow in Non-Autonomous Quadratic Maps and the Riemann Zeros: A Numerical Study,''} for consideration as a research article in \textit{Fractal and Fractional}.

\vspace{0.1em}

\noindent\textbf{Significance and context.} Finding a dynamical system whose spectrum reproduces the non-trivial Riemann zeros is a long-standing question in mathematical physics (the Hilbert--P\'olya conjecture). Existing analytical approaches, from Berry--Keating's $H=xp$ Hamiltonian to random-matrix models, elegantly capture the macroscopic statistics of high-order zeros but rarely generate the discrete low-order zeros deterministically. We take a different, ``bottom-up'' route: building on our recently published result that the symbolic dynamics of the Logistic map at its band-merging point is isomorphic to the prime sieve, we introduce a non-autonomous quadratic map governed by a logarithmic renormalization-flow schedule and study its transfer-operator spectrum. The model reproduces the low-order Riemann zeros with small numerical error and exhibits a change of behavior near $N\approx20$ that shows a qualitative spatial coincidence with the enlarged error bars reported in recent USTC ion-trap quantum simulations of the same zeros. At large $N$, an ablation study shows that restoring the conjugate spectrum drives the fitted ground-state intercept to zero, pointing to single-sided truncation as a likely source of the residual envelope.

\vspace{0.1em}

\noindent\textbf{Why \textit{Fractal and Fractional}.} The manuscript centers on renormalization-flow scaling, non-autonomous (time-dependent) discrete dynamics, and the fractal/self-similar structure of the logistic map's band-merging cascade, connecting nonlinear dynamical-systems theory to a classical problem in analytic number theory. This combination of fractal dynamics, scaling laws, and numerical spectral analysis fits well within the journal's scope.

\vspace{0.1em}

\noindent\textbf{On rigor and claims.} We are deliberately explicit about epistemic status throughout the manuscript: results are presented as numerical observations and heuristic arguments, not proofs, and a dedicated table classifies every main claim as established, published, numerical, heuristic, qualitative, or conjectural. We do not claim to have identified the Hilbert--P\'olya operator or to have proven any statement about the Riemann Hypothesis. All code and data are openly available at \url{https://github.com/maris205/riemann_logistic} for full reproducibility.

\vspace{0.1em}

\noindent\textbf{Prior submissions.} This manuscript has not been previously submitted to any MDPI journal.

\vspace{0.1em}

\noindent We confirm that neither the manuscript nor any parts of its content are currently under consideration for publication with or published in another journal.

\vspace{0.1em}

\noindent All authors have approved the manuscript and agree with its submission to \textit{Fractal Fract.}

\vspace{0.15em}

\noindent Thank you for considering our submission. We would be happy to suggest suitable referees with expertise in nonlinear dynamics, ergodic theory, and analytic number theory upon request.

\vspace{0.1em}

\noindent Sincerely,

\vspace{0.15em}

\noindent Liang Wang \\
School of Artificial Intelligence and Automation, \\
Huazhong University of Science and Technology, Wuhan, P.R. China \\
\href{mailto:wangliang.f@gmail.com}{wangliang.f@gmail.com}

\end{document}
